\section{Can the constant be improved?}\label{ch18:sec:sharp}
Our proof used the triangle inequality, which may give away more than necessary. Could a cleverer proof give a smaller factor than $1/(1-q)$ in part~2 of Theorem~\ref{thm:inexact}? To find out, we test the bound on the assistant's example: $X=\RR$, $T(x)=qx$, $y_0=0$, a number $\eta\ge0$, and
\[
 y_{n+1}=qy_n+\eta\qquad\text{for every }n\in\NN.
\]
As Section~\ref{ch18:sec:guess} checked, $T$ is a contraction with factor $q$ and fixed point $x_*=0$, and every update error is $\eta$. The first points $y_1=\eta$, $y_2=q\eta+\eta$ and $y_3=q^2\eta+q\eta+\eta$ suggest
\[
 y_n=\eta\sum_{k=0}^{n-1}q^k
     =\eta\,\frac{1-q^n}{1-q}\qquad\text{for every }n\in\NN.
\]
A pattern is only a guess, so we prove the formula by induction. At $n=0$ it gives $0=y_0$. If it holds for some $n$, then
\begin{align*}
 y_{n+1}
 &=qy_n+\eta\\
 &=q\,\frac{\eta(1-q^n)}{1-q}+\eta\\
 &=\frac{\eta(q-q^{n+1})+\eta(1-q)}{1-q}\\
 &=\frac{\eta(1-q^{n+1})}{1-q},
\end{align*}
which is the formula for $n+1$. The induction never used the sign of $\eta$, so the formula holds for every real number $\eta$, negative ones included. Since $y_n-\eta/(1-q)=-\eta q^n/(1-q)$ and $q^n\to0$, the points approach $\eta/(1-q)$, as the assistant claimed. This is no accident: the biased map $y\mapsto qy+\eta$ is itself a contraction, and its fixed point, the solution of $y=qy+\eta$, is $\eta/(1-q)$, the limit that Theorem~\ref{thm:banach} predicts.

Now compare with the theorem. Here $e_0=0$, and since $\eta\ge0$ and $0<q^n\le1$, every $y_n$ is nonnegative, so
\[
 e_n=y_n=\eta\,\frac{1-q^n}{1-q}.
\]
This is exactly the bound of part~2, $q^n\cdot0+\eta(1-q^n)/(1-q)$: the bound holds with equality at every step. So the factor $1/(1-q)$ cannot be replaced by a smaller number. Suppose someone claims that for some $K<1/(1-q)$, every iteration satisfying our assumptions, with all update errors at most $\eta$, obeys $e_n\le q^ne_0+K\eta$ for every $n\ge1$. In the example with $\eta>0$, the claim says $\eta(1-q^n)/(1-q)\le K\eta$. Dividing by $\eta$ and multiplying by $1-q$, both positive, gives $1-q^n\le K(1-q)$, that is,
\[
 q^n\ge1-K(1-q).
\]
The right side is positive, because $K(1-q)<1$, and $q^n\to0$, so the inequality fails for every sufficiently large $n$.

An upper bound that holds for every example in some class is called \emph{sharp} if it cannot be lowered: every smaller proposed bound fails for some example in the class. So the factor $1/(1-q)$ in part~2 is sharp for the class of all iterations that satisfy our assumptions.

A sharp bound on the distance to $x_*$ still does not tell us whether the points converge. Let $T(x)=0$ on $\RR$, a contraction with every factor $q$ in $(0,1)$ and fixed point $0$. Let $\eta>0$ and $y_n=(-1)^n\eta$. Every update error is $|(-1)^{n+1}\eta-0|=\eta$, and every error $e_n$ equals $\eta$, within the bound of part~2. Yet the points jump back and forth between $\eta$ and $-\eta$ and do not converge (Exercise~\ref{ex:18.3}). That is why part~3 needs the stronger hypothesis $\eps_n\to0$.
